\documentclass[12pt]{article}
\usepackage[margin=0.5in]{geometry}
\usepackage{enumitem}
\usepackage{graphicx}

\title{USC Facilities Rain Barrel}
\author{EMCH 427-002 Group 23}
\date{}

\begin{document}

\maketitle

\section*{Project Overview}
The Green Quad, the first LEED-certified building on campus, is equipped with a 1000-gallon rainwater tank connected to the building's downspouts. The current method of manually transferring water using 2-gallon cans is time-consuming and inefficient. This project aims to design a system to transport water from the tank to the garden in a time-efficient, environmentally friendly, and cost-effective manner, adhering to the client's specifications.

\section*{Client Requirements}
\begin{itemize}[noitemsep, topsep=0pt]
\item Schedule watering on Monday, Wednesday, Thursday, and twice on Friday.
\item Ensure no new environmental impact.
\item Option to add a permanent fixture to the tank.
\item High water pressure suitable for sprinklers.
\item Simplicity for operation by untrained individuals.
\item Low maintenance, with acceptable yearly maintenance.
\item Operation by a single person preferred.
\item Comparable to city water flow rate and pressure.
\item Long-term usability.
\item Minimal aesthetic concerns; not a major eyesore.
\item Noise level not exceeding that of a battery-powered leaf blower.
\item Budget: Preferably under \$1,000, maximum \$2,000.
\item Preference for reusing materials and eco-friendly solutions.
\item Prevent clogging from roof debris.
\item Fill a 2-gallon can in 30 seconds, preferably in 10 seconds.
\item Trees and other equipment can be relocated if necessary.
\end{itemize}

\newpage
\section*{Proposed Solutions}

\subsection*{1. Gravity-Fed System}

\noindent\textit{Description:} Utilize gravity to create sufficient water pressure by elevating the tank.

\noindent\textit{Implementation Steps:}
\begin{enumerate}[noitemsep, topsep=0pt]
\item \textbf{Increase Tank Height:} Raise the elevation of the existing tank to increase the water pressure. Since the tank is currently 5 feet above the garden but effectively at -2 feet due to the loading dock position, repositioning is necessary.
\item \textbf{Structural Support:} Install a sturdy support structure to elevate the tank safely above the required height.
\item \textbf{Piping Adjustments:} Replace the current 2-inch threaded pipe outlet with a larger diameter pipe if necessary to reduce friction losses and maintain flow rate.
\item \textbf{Pressure Regulation:} Incorporate a pressure regulator to ensure consistent water pressure for sprinklers.
\item \textbf{Filtration System:} Install a debris filter at the outlet to prevent clogging from roof runoff.
\end{enumerate}

\noindent\textit{Advantages:}
\begin{itemize}[noitemsep, topsep=0pt]
\item No additional energy consumption.
\item Low maintenance once installed.
\item Simple operation requiring minimal training.
\end{itemize}

\noindent\textit{Disadvantages:}
\begin{itemize}[noitemsep, topsep=0pt]
\item Structural modifications may increase costs.
\item Limited by the achievable height based on existing infrastructure.
\end{itemize}
\newpage
\subsection*{2. Manual Pump System}

\noindent\textit{Description:} Employ a manual pump, such as a treadle or hand pump, to move water from the tank to the garden.

\noindent\textit{Implementation Steps:}
\begin{enumerate}[noitemsep, topsep=0pt]
\item \textbf{Pump Selection:} Choose a durable manual pump capable of delivering the required flow rate (2 gallons in 30 seconds).
\item \textbf{Mounting Fixture:} Install a permanent mounting point on the tank for the manual pump.
\item \textbf{Piping and Hose Setup:} Connect hoses from the pump to the garden area, ensuring minimal length to reduce effort.
\item \textbf{Handle Design:} Design ergonomic handles for ease of use by a single operator.
\item \textbf{Debris Filter:} Integrate a filter to prevent clogging from roof debris.
\end{enumerate}

\noindent\textit{Advantages:}
\begin{itemize}[noitemsep, topsep=0pt]
\item Low initial cost.
\item No electricity required.
\item Eco-friendly with minimal environmental impact.
\end{itemize}

\noindent\textit{Disadvantages:}
\begin{itemize}[noitemsep, topsep=0pt]
\item Physical effort required may be tiring.
\item Achieving high pressure comparable to city water may be challenging.
\item Potentially slower filling times without mechanical advantage.
\end{itemize}
\newpage
\subsection*{3. Powered Pump System}

\noindent\textit{Description:} Utilize an electric or solar-powered pump to automate water transfer, ensuring high pressure and flow rate.

\noindent\textit{Implementation Steps:}
\begin{enumerate}[noitemsep, topsep=0pt]
\item \textbf{Pump Selection:} Select a submersible or external electric pump compatible with the tank’s outlet and garden hose.
\item \textbf{Power Source Setup:} Install the pump near a reliable power source or incorporate solar panels for sustainable energy.
\item \textbf{Pressure Booster:} Integrate a pressure booster to achieve the desired high-pressure output for sprinklers.
\item \textbf{Control System:} Implement simple controls (e.g., switches or timers) for ease of operation by untrained personnel.
\item \textbf{Filtration:} Add a debris filter to protect the pump and prevent clogging.
\item \textbf{Safety Features:} Incorporate automatic shutoff mechanisms to prevent dry running and overheating.
\end{enumerate}

\noindent\textit{Advantages:}
\begin{itemize}[noitemsep, topsep=0pt]
\item High flow rate and pressure suitable for sprinklers.
\item Minimal physical effort required.
\item Automated operation can save time.
\end{itemize}

\noindent\textit{Disadvantages:}
\begin{itemize}[noitemsep, topsep=0pt]
\item Higher initial cost, potentially approaching the budget limit.
\item Requires access to electricity or investment in solar infrastructure.
\item Increased maintenance compared to gravity-fed systems.
\end{itemize}
\newpage
\subsection*{4. Air-Pressure System}

\noindent\textit{Description:} Use compressed air to create pressure for water distribution, either through a manual pump or a hand-powered mechanism.

\noindent\textit{Implementation Steps:}
\begin{enumerate}[noitemsep, topsep=0pt]
\item \textbf{Air Compressor Selection:} Choose a suitable manual or hand-powered air compressor capable of generating the necessary pressure.
\item \textbf{Integration with Tank:} Connect the air compressor to the tank’s outlet, ensuring airtight seals.
\item \textbf{Pressure Tank Installation:} Install a pressure tank to store compressed air for consistent water flow.
\item \textbf{Regulation Controls:} Implement pressure regulators to maintain steady water pressure for sprinklers.
\item \textbf{Hose and Sprinkler Setup:} Connect hoses from the pressure system to the garden sprinklers.
\item \textbf{Debris Filtration:} Add filters to prevent clogging from roof debris.
\end{enumerate}

\noindent\textit{Advantages:}
\begin{itemize}[noitemsep, topsep=0pt]
\item Can achieve high pressure without electricity.
\item Potential for fast water transfer once pressure is established.
\item Environmentally friendly with reusable components.
\end{itemize}

\noindent\textit{Disadvantages:}
\begin{itemize}[noitemsep, topsep=0pt]
\item Initial setup complexity may be higher.
\item Manual operation of the air compressor can be labor-intensive.
\item Maintenance of air seals and compressor components required.
\end{itemize}
\newpage
\subsection*{5. Hybrid System}

\noindent\textit{Description:} Combine elements from multiple systems to optimize performance, reliability, and ease of use.

\noindent\textit{Implementation Steps:}
\begin{enumerate}[noitemsep, topsep=0pt]
\item \textbf{Gravity Enhancement:} Slightly elevate the tank to utilize residual gravity force.
\item \textbf{Solar-Powered Pump Integration:} Incorporate a solar-powered pump to supplement gravity, ensuring consistent pressure and flow.
\item \textbf{Automated Controls:} Use simple timers or sensors to automate the watering schedule.
\item \textbf{Comprehensive Filtration:} Implement multi-stage filters to handle debris and prevent clogging.
\item \textbf{User-Friendly Interface:} Design intuitive controls for ease of operation by a single individual.
\end{enumerate}

\noindent\textit{Advantages:}
\begin{itemize}[noitemsep, topsep=0pt]
\item Balances energy efficiency with automation.
\item Redundancy increases system reliability.
\item Flexible adaptation to varying conditions and usage.
\end{itemize}

\noindent\textit{Disadvantages:}
\begin{itemize}[noitemsep, topsep=0pt]
\item Potentially higher initial costs.
\item More complex system requiring careful integration.
\item Maintenance may be more involved due to multiple components.
\end{itemize}
\newpage
\subsection*{6. Drip Irrigation System}

\noindent\textit{Description:} Implement a drip irrigation system to deliver water directly to the plant roots efficiently, reducing water usage and ensuring consistent moisture levels.

\noindent\textit{Implementation Steps:}
\begin{enumerate}[noitemsep, topsep=0pt]
\item \textbf{System Design:} Layout the drip lines to cover the garden area, ensuring even distribution of water.
\item \textbf{Pump Integration (Optional):} If using a powered pump, integrate it to provide consistent pressure for the drip emitters.
\item \textbf{Filtration:} Install fine filters to prevent clogging of drip emitters from roof debris.
\item \textbf{Control Valves:} Incorporate timers or manual valves to automate the watering schedule.
\item \textbf{Hose Connections:} Connect the drip lines to the tank outlet or pump output using appropriate fittings.
\item \textbf{Maintenance Access:} Design the system for easy access to emitters and filters for regular maintenance.
\end{enumerate}

\noindent\textit{Advantages:}
\begin{itemize}[noitemsep, topsep=0pt]
\item Highly efficient water usage with minimal waste.
\item Reduced evaporation and runoff compared to sprinklers.
\item Can be easily automated with timers for scheduled watering.
\end{itemize}

\noindent\textit{Disadvantages:}
\begin{itemize}[noitemsep, topsep=0pt]
\item Initial setup can be labor-intensive.
\item Emitters may clog if filtration is inadequate.
\item Limited coverage area compared to sprinkler systems unless extensively installed.
\end{itemize}

\end{document}
